\section{Introduction}
    \IEEEPARstart{I}{\lowercase{nterferometric \uppercase{S}ynthetic \uppercase{A}perture \uppercase{R}adar}} (\acs{InSAR}) analysis, a pivotal technique in 
    %\textcolor{blue}{\acl{SAR} (\acs{SAR})} 
    \acl{SAR} (\acs{SAR}) imagery, is widely employed for various applications such as  monitoring volcanoes and earthquakes, assisting in construction and urban planning, investigating the effects of climate changes, among others. This technique operates by transmitting  radar waves from the satellite to the ground and capturing the backreflected echoes. Ground movements with high accuracy can be detected by comparing phases of \acs{SAR} images acquired at two different dates (2-pass \acs{InSAR}). However, this technique shows limitations when image acquisitions become distant, which results in coherence loss, thus information loss. %\textcolor{blue}{\acl{MT-InSAR} (\acs{MT-InSAR})}
    \acl{MT-InSAR} (\acs{MT-InSAR}) approaches offer enhanced performances compared to the $2$-pass approach. One of the most widely used approaches in \acs{MT-InSAR} is %\textcolor{blue}{\acl{PL} (\acs{PL})}
    \acl{PL} (\acs{PL}), which exploits the full covariance matrix of 
    %\textcolor{blue}{ \acl{SARTS} (\acs{SARTS})} 
    \acl{SARTS} (\acs{SARTS}) and their statistical properties. %\yysup{These methods make use of the full coherence matrix in a mathematical framework, ensuring the robustness against phase biases caused by short-lived fading signals in short temporal baseline interferograms \citep{ansari2020study}.} 
    The first objective of \acs{PL} approaches consists of estimating the phase differences relative to a reference image from the \acs{SARTS}, with improved accuracy and respecting the temporal phase closure. For this, early PL approaches require prior knowledge on coherence matrix and take the modulus of the \ac{SCM} as a plug-in for the coherence matrix \citep{guarnieri2008exploitation}. This approach is equivalent to taking the 
    %\textcolor{blue}{\acl{MLE} (\ac{MLE})} 
    \acl{MLE} (\ac{MLE}) of the covariance matrix without considering its structure (see equation \eqref{Equation:cov_mat_struc}) which makes it a suboptimal estimator due to the use of the modulus operator which is non-holomorphic. 
    %A proposed solution \citep{vu2022new, 9763551} assumes a Gaussian model and jointly estimates the coherence matrix along with the phases, demonstrating improvements on simulations and real data. 
    %\yy{proposed an improved solution in which phase differences and coherence matrix are jointly estimated, assuming both Gaussian and generalized gaussian models.} \yysup{
    To correct this issue, a true \ac{MLE} framework has been proposed in \citep{vu2022new}, where the coherence and the phases are estimated by a $1$-step approach. The proposed solution involves using iterative algorithms such as \ac{BCD} including a \ac{MM}. % to estimate the coherence matrix and the phases vector. 
    An extension to non Gaussian data has been proposed \citep{vu2023robust}, which shows their interest in urban areas where the Gaussian assumption is often no longer valid \citep{mian2018new}.
    %\textcolor{blue}{
    However, this algorithm remains computationally intensive, particularly due to use of iterative procedures and complex mathematical operations on matrices that grow in size over time. %}
    %\textcolor{blue}{
    Additionally, it is unable to efficiently handle the integration of new images and requires a complete re-estimation of all phases of the time series as well as the covariance matrix, which in this case corresponds to the Tyler estimator, which itself is computationally expensive to calculate. This becomes a real issue especially with the vast amount of available \acs{SAR} images thanks to ongoing \acs{SAR} missions.%}
    
    \noindent
    %\textcolor{blue}{ 
    Recently, online signal processing algorithms have seen a great attraction and interest due to their ability to process data continuously in real time, unlike batch algorithms that require additional time to gather sufficient samples \citep{zakharov2008low, angelosante2010online, elvander2016online, zhang2023high}. This real-time ability eliminates the need for complex and resource-intensive operations, making online approaches a more efficient and cost-effective solution. %}
    
    \noindent
    %\textcolor{blue}{
    To the best of our knowledge, sequential processing of \acs{SARTS} has not been extensively explored. One straightforward approach is to consider phase differences only between the new image and several previous images with close acquisition dates, like the \ac{SBAS} method \citep{lanari2004small} which pairs \ac{SAR} images with small temporal baselines.
    However, this approach can introduce phase bias due to the presence of short-lived fading signals \citep{ansari2020study}.
    One notable sequential approach was proposed in \citep{ansari2017sequential}, where the idea is to divide the \acs{SARTS} into mini-stacks and apply the classic \acs{PL} to each mini-stack. The mini-stacks are then compressed into a single virtual image through \ac{PCA}. Each mini-stack is treated relative to its first image, a global \acs{PL} is performed to link all the mini-stacks through the compressed version of each mini-stack. %\textcolor{blue}{The drawbacks of this approach include using classic \acs{PL} within a Gaussian model framework and determining an appropriate  period for \acs{SAR} image acquisition.} 
    This approach use \acs{PL} within a Gaussian framework and necessites the determination of an appropriate period for sufficient \acs{SAR} image acquisition.
    In this paper, we propose a sequential approach based on \citep{vu2023robust}, which leverages the benefits of the \acs{MLE} over classic \acs{PL} while maintaining a lower computational cost. It generalizes the approach proposed in \citep{elhajjar2024}, namely \ac{SgMLEPL}, to a Non-Gaussian model. 
    %In \citep{elhajjar2024}, another sequential approach is proposed to integrate one image at once, based on a maximum likelihood estimator in a Gaussian framework. The key idea is to partition the data vector into two blocks: one representing previous data and another representing the new image.
    %}
    
    %\citep{mian2018new, vu2023robust} showed that \acs{SAR} data fits a generalized Gaussian model rather than a purely Gaussian one in urban areas.
    
    %Several \acs{PL} techniques has been already proposed, such as SqueeSAR \citep{ferretti2011new}, \ac{EMI} \citep{ansari2018efficient}, \ac{MLEPL} \citep{vu2022new, vu2023robust, 9763551}. 
    %By repeating this process over multiple satellites passes (e.g., Sentinel-1 mission $6-12$ days), sequences of \acs{SAR} images can be generated. 
    %\acs{MT-InSAR} approaches offers enhanced performances compared with the $2$-pass approach, which consists in computing the phase differences using $2$ \acs{SAR} images.  
    %These methods exploit the coherence matrix (i.e target statistics) in a mathematical framework, ensuring the robustness against phase biases caused by short-lived fading signals in short temporal baseline interferograms \citep{ansari2020study}. The classic \acs{PL} method estimated phase differences based of the time series of input \acs{SAR} images, with the coherence matrix typically chosen as the modulus of the \ac{SCM}. However, this approach is equivalent to taking the \acs{MLE} of the covariance matrix without considering its structure (\ref{Equation:cov_mat_struc}) which make it a suboptimal estimator due to the non-holomorphic nature of the $|.|$ operator. A proposed solution in \citep{vu2022new} assumes a Gaussian model and jointly estimates the coherence matrix along with the phases, demonstrating improvements on simulations and real data. Moreover, \citep{vu2023robust} showed that \acs{SAR} images fits a generalized Gaussian model rather than a purely Gaussian one.
    %Furthermore, it is worth mentioning that ongoing \acs{SAR} missions regularly provide a vast amount of \acs{SAR} images. As a consequence, applying offline \acs{PL} approaches each time when a newly acquired image arrives involves substantial computational operations and require immense computation time due to the increase in data vector size with the addition of a new component. % \yy{The integration of new images is thus not efficient with the offline PL approaches. }
    %In \citep{vu2023robust}, the proposed solution involves using iterative algorithms such as \ac{BCD} including a \ac{MM} to estimate the coherence matrix and the phases vector. This iterative process effectively re-estimates the covariance matrix as it grows. 
    %To the best of our knowledge, sequential processing of \acs{SARTS} has not been extensively explored. One straightforward approach is to consider phase differences only between the new image and several previous images with close acquisition dates, like the \ac{SBAS} method \citep{lanari2004small} which pairs \ac{SAR} images with small temporal baselines. % such as the \acs{SBAS} method which uses small spatial and temporal baselines between \acs{SAR} image pairs. This approach introduces a bias 
    %However, this approach can introduce phase bias due to the presence of short-lived fading signals \citep{ansari2020study}.
    %One notable sequential approach was proposed in \citep{ansari2017sequential}, where the idea is to divide the \acs{SARTS} into mini-stacks and apply the classic \acs{PL} to each mini-stack. The mini-stacks are then compressed into a single virtual image through \ac{PCA}. Each mini-stack is treated relative to its first image, a global \acs{PL} is performed to link all the mini-stacks through the compressed version of each mini-stack. %\textcolor{blue}{The drawbacks of this approach include using classic \acs{PL} within a Gaussian model framework and determining an appropriate  period for \acs{SAR} image acquisition.} 
    %\textcolor{red}{This approach use classic PL within a Gaussian framework and necessites the determination of an appropriate period for sufficient SAR image acquisition.} \textcolor{red}{COMMENT: je ne pense pas que c'est une faiblesse d'utiliser PL classique et modèle gaussian, ce dernier est valide dans beaucoup de cas.}
    %This approach use classic \acs{PL} within a Gaussian framework and necessites the determination of an appropriate period for sufficient \acs{SAR} image acquisition.
    %In \citep{elhajjar2024}, another sequential approach is proposed to integrate one image at once, based on a maximum likelihood estimator in a Gaussian framework. The key idea is to partition the data vector into two blocks: one representing previous data and another representing the new image.
    
    \noindent
    In this paper, we present a robust sequential approach. %% based on \citep{elhajjar2024, vu2023robust}. %\yy{In this approach, conditional distribution of the new image on previous images is estimated ... }
    Specifically, the previous images along with their corresponding estimated parameters are used as prior information, allowing us to leverage the conditional distribution of the new image relative to the previous ones. Our approach resembles those used in handling missing data, where prior information plays a crucial role \citep{aubry2021structured, hippert2022robust}.
    %\textcolor{red}{A détailler plus la méthodologie, nouvelle cov conditionnée sur cov past, param à estimer}
    Closed-form solutions for each parameter are derived and computed using a \acs{BCD} algorithm. Our approach demonstrates similar performance to traditional offline methods in both simulations and real data studies, while also significantly reducing computational costs. %Additionally, it offers the advantage of reduced computational costs compared to offline approaches.
    %The rest of the paper is prganized as folows: section \ref{section:PL} presents the \acs{PL} problem. Section \ref{section:data_model} introduces our data model and the covariance matrix structure, followed by our contribution in section \ref{section:approach}. Validation on simulations and real data are provided in sections \ref{section:simu} and \ref{section:real_data}.
    
    \noindent
    %\textcolor{blue}{
    In the continuation of the paper, the exponent $^T$ (resp. $^*$) denotes the transposed operator (resp. the conjugated operator), while the symbol $\circ$ characterizes the element-wise (Hadamard) multiplication, and the exponent $^H$ corresponds to the transposed and complex conjugated operator (Hermitian). Matrices are written in bold uppercase, vectors in bold lowercase, and scalars in lowercase.%}
    %, and the exponent $^H$ corresponds to the transposed and complex conjugated operator (Hermitian). 
    